\documentclass[10pt,a4paper]{article}
\usepackage{amsmath,amssymb,amsthm,graphicx}
\usepackage{amsmath,amssymb,amsthm}
\usepackage{hyperref}
\newtheorem{theorem}{Theorem}
\newtheorem{definition}[theorem]{Definition}
\newtheorem{remark}[theorem]{Remark}
\newtheorem{proposition}[theorem]{Proposition}
\newtheorem{corollary}[theorem]{Corollary}
\newcommand{\C}{\mathbb C}
\newcommand{\Q}{\mathbb Q}
\newcommand{\N}{\mathbb N}
\newcommand{\T}{\operatorname T}
\newcommand{\ST}{\operatorname{ST}}
\newcommand{\SL}{\operatorname{SL}}
\newcommand{\GL}{\operatorname{GL}}
\newcommand{\Sym}{\operatorname{Sym}}
\newcommand{\cC}{\mathcal C}
\newcommand{\cN}{\mathcal N}
\newcommand{\cM}{\mathcal M}
\newcommand{\ba}{\textbf a}
\newcommand{\bb}{\textbf b}
\newcommand{\bc}{\textbf c}
\newcommand{\bbm}{\textbf m}
\newcommand{\bx}{\textbf x}
\newcommand{\by}{\textbf y}
\newcommand{\bz}{\textbf z}
\newcommand{\bo}{\textbf 0}
\newcommand{\poly}{\operatorname{poly}}
\newcommand{\tensor}{\textstyle\bigotimes}
\DeclareMathOperator{\SU}{SU}
\DeclareMathOperator{\Pf}{Pf}
\DeclareMathOperator{\per}{Per}
\DeclareMathOperator{\sgn}{sgn}
\newcommand{\la}{\lambda}
\newcommand{\VP}{\textup{VP}}
\newcommand{\VNP}{\textup{VNP}}
\AtBeginDocument{\author{Ankit Garg, Christian Ikenmeyer, Visu Makam\\
Rafael Oliveira, Michael Walter, Avi Wigderson}}
\title{Succinct generators for SL tensor actions would collapse VP and VNP through hyperpfaffian completeness}
\author{Ankit Garg, Christian Ikenmeyer, Visu Makam, Rafael Oliveira, Michael Walter, Avi Wigderson}
\date{doi:10.4230/LIPIcs.CCC.2020.12; CC BY 3.0 Unported}
\begin{document}
\maketitle
\noindent\textit{Editorial note.} Editable transcription only of formally published six-author CC BY3.0 Unported content. CompleteSection4 retained, including all author source-local representation analysis, hyperpfaffian completeness and final encoding deduction. CompleteSection3 torus proof is shared uncounted context because the final author paragraph explicitly invokes its deduction; no replacement proof supplied. Nativev2 extra tensor-action section excluded and raw files private. Original author equations/indices/assumptions retained; no mathematical correction or new assumption. OriginalSection1.2/Conjecture6 external references presented accurately without importing framework as a new result. All production/validation onlyYupeng; root owns compilation.\par\medskip
\noindent\textit{Editorial note.} Revisionr002: split only the six-author title block into two lines via a begin-document author hook; the metadata authors array is unchanged. The three actual outputpage8 overflowing displays are rewrapped between tensor factors, preserving every original summand, factor, index, exponent, parenthesis and equality.\par\medskip
\noindent\textit{Allowed standard prerequisites (editorial).}\begin{itemize}
\item Invariant theory ofSL representations, Schur-Weyl duality and algebraic circuit complexity VP/VNP/projections. EntireformalSection4 representation/symmetric-tensor definitions and standard SchurWeyl statement retained; source-local uniqueness/hyperpfaffian projection proof fully included, Section2/3 deduction context retained
\end{itemize}\medskip
\section*{Group actions and invariant generators (original Section1)}
An action (also called a representation) of a group $G\subseteq\GL_n(\C)$ on an $m$-dimensional complex vector space $V$ is a group homomorphism $\pi:G\to\GL_m(\C)$, that is, an association of an invertible $m\times m$ matrix $\pi(g)$ for every group element $g\in G$, satisfying $\pi(g_1g_2)=\pi(g_1)\pi(g_2)$ for all $g_1,g_2\in G$ (and $\pi(e)=I_m$, where $e\in G$ is the identity element and $I_m$ is the identity matrix). A group element $g\in G$ acts on a vector $v\in V$ by the linear transformation $\pi(g)$, and in this paper we will be dealing with algebraic actions, that is, the entries of the matrix $\pi(g)$ will be rational functions in the entries of the matrix $g$. We will write $g\cdot v=\pi(g)v$. We denote by $\C[V]$ the ring of polynomial functions on $V$.
A polynomial function $f\in\C[V]$ is called invariant if $f(g\cdot v)=f(v)$ for all $v\in V$ and $g\in G$. The induced action on polynomials is $(g\cdot p)(v)=p(g^{-1}\cdot v)$. A set $\{f_i\}_{i\in I}$ of invariant polynomials is called a generating set if any other invariant polynomial can be written as a polynomial in the $f_i$'s. The invariant ring is denoted $\C[V]^G$.
\setcounter{theorem}{2}
\begin{definition}[\textbf{Succinct encoding of generators}] \label{def:succ.enc}
	Fix an action of a group $G$ on a vector space $V = \C^m$. We say that an arithmetic circuit
	$\cC(x_1,\ldots, x_m,y_1,\ldots, y_r)$ \emph{succinctly encodes} the generators of the invariant
	ring if the set of polynomials formed by evaluating the $y$-variables,
	$\{\cC(x_1,\ldots, x_m, \alpha_1,\ldots, \alpha_r)\}_{\alpha_1,\ldots, \alpha_r \in \C}$,
  is a generating set for the invariant ring~$\C[V]^G$.
\end{definition}

\begin{remark}
	The \emph{size} of a succinct encoding as defined above is given by the size of the circuit
	$\cC(x_1, \ldots, x_m, y_1, \ldots, y_r)$, which is measured by the bit complexity of the
	constants used in the computation of $\cC$ as well as the number of gates of the computation
	graph of $\cC$.
  In particular, this means that all constants used in the computation of $\cC$ are rationals.
\end{remark}

\setcounter{theorem}{10}
\begin{theorem}\label{MC-cex}
Let $k \geq 2$ be even.
Consider the action of $G = \SL_{2kn}(\C)$ on $V = \tensor^{2k} \C^{2kn}$. Then any set of generators for the invariant ring cannot have a polynomial sized (in $n$) succinct encoding, unless ${\rm VP = VNP}$.
\end{theorem}

\setcounter{theorem}{7}
For our first counterexample, we analyze a simple (torus) action on $3$-tensors.
Here, $\ST_n(\C)$ denotes the group of $n\times n$ diagonal matrices with determinant $1$.

\begin{theorem}\label{thm:torus-action}
Consider the natural action of $G = \ST_n(\C) \times \ST_n(\C) \times \ST_n(\C)$ on $V = \C^n \otimes \C^n \otimes \C^n$. Then any set of generators for the invariant ring cannot have a polynomial sized (in $n$) succinct encoding, unless ${\rm NP \subseteq P/poly}$.
\end{theorem}

\setcounter{section}{1}
\section{Preliminaries}
\setcounter{theorem}{15}
%

In this section we establish notation and we formally state basic facts and definitions which
we will need in later sections.



\begin{definition}[3-dimensional matching~\cite{K72}]\label{def:3D-matching}
	The \emph{3-dimensional matching} problem is defined as follows:
	\begin{itemize}
		\item[Input:] a set $U \subseteq [n] \times [n] \times [n]$, representing the edges of
		a tripartite, 3-uniform hypergraph.
		\item[Output:] YES, if there is a set of hyperedges $W \subseteq U$ such that $|W| = n$ and
		no two elements of $W$ agree in any coordinate (that is, they form a matching in this
		hypergraph). NO, if there is no such set.
	\end{itemize}
\end{definition}

\begin{theorem}[NP-completeness of 3-dimensional matching~\cite{K72}]\label{thm:3D-matching-NP-complete}
	The 3-dimensional matching problem is NP-complete.
\end{theorem}

\subsection{Basic facts from algebraic complexity}

We now give basic facts that from algebraic complexity which we will use in the next sections.

The next proposition shows that homogeneous components of low degree
of an arithmetic circuit can be efficiently computed, with a small blow-up
in circuit size and without the use of any extra constants. This proposition was
originally proved by Strassen in~\cite{S73} and its proof can be found
in~\cite[Theorem 2.2]{SY10}.
In the following proposition, given a polynomial $p(\bx)$, we denote its degree-$d$ homogeneous component by $H_d[p(\bx)]$.

\begin{proposition}[Efficient computation of homogeneous components]\label{prop:hom-components}
	Given a circuit $\cC(\bx)$ of size $s$, then for every $r \in \N$
	there is a homogeneous circuit $\Psi(\bx)$ of size $O(r^2 s)$
	computing $H_0[\cC(\bx)], H_1[\cC(\bx)], \ldots, H_r[\cC(\bx)]$.
	Moreover, the constants used in the computation of the components $H_i[\cC(\bx)]$
	are a subset of the coefficients used in the computation of $\cC(\bx)$.
%
%
\end{proposition}

The next theorem, proved by~\cite[Theorem 4.10]{AB03} gives us a randomized polynomial time
algorithm to test whether an algebraic circuit of polynomial size, with rational
coefficients, is identically zero. Another randomized algorithm easily follows
from~\cite[Lemma 2]{S79}, when adapted for polynomials with rational coefficients.

\begin{theorem}[PIT for poly-sized circuits~\cite{AB03}]\label{thm:PIT-BPP}
	Let $P(\bx) \in \Q[\bx]$ be a polynomial in the variables $\bx = (x_1, \ldots, x_n)$, with
	each variable $x_i$ having degree bounded by $d_i$, and whose coefficients are rationals with
	bit complexity bounded by $B$. If $P(\bx)$ is given as an arithmetic circuit of size $s$, then
	there exists a randomized algorithm running in time ${\rm poly}(n, s, \log(B), 1/\epsilon)$
	and using $O\left(\sum_{i=1}^n \log(d_i) + \log(B) \right)$ random bits which tests whether $P(\bx)$ is identically zero.
	If $P(\bx)$ is the identically zero polynomial then the algorithm always succeeds.
	Otherwise, it errs with probability at most $\epsilon$.
\end{theorem}

\setcounter{theorem}{19}
\section{Hardness of Generators for torus actions}\label{sec:torus}

Let $\C^*$ denote the multiplicative group consisting of all non-zero
complex numbers. A direct product $\T_n = (\C^*)^n$ is called a torus, and
is clearly an abelian group. Tori are important examples of reductive
groups -- any abelian connected reductive group is a torus! It is often the case that it is easier to understand tori in comparison with more general (non-abelian) reductive groups. This is no different for invariant theory, see for example \cite{DK,Wehlau}. We also point to \cite[Proposition~3.3]{DM-exp} for an elementary linear algebraic description of the invariant ring for torus actions. Conjecture 6 of the original publication already fails in this well behaved setting.
This is the content of our Theorem~\ref{thm:torus-action}, which we will prove in this section.
Recall that $\ST_n(\C) \cong \{ \bz \in \T_n : z_1\cdots z_n = 1 \}$, which is itself a torus.

\begin{theorem}[Theorem~\ref{thm:torus-action}, restated]
	Consider the natural action of $G = \ST_n(\C) \times \ST_n(\C) \times \ST_n(\C)$
	on $V = \C^n \otimes \C^n \otimes \C^n$, where an element
	$(\ba, \bb, \bc) \in G$ acts on a tensor $u \in V$ as follows:
  $(\ba,\bb,\bc) \cdot u := v$, such that
	$v_{ijk} = a_i b_j c_k u_{ijk}$. Any set of generators
	for the invariant ring of this action cannot have a polynomial sized
	(in $n$) succinct encoding, unless ${\rm NP \subseteq P/poly}$.
\end{theorem}

\begin{proof}
	Suppose that the natural action above has a set of generators with
	a polynomial sized succinct encoding. Thus, there is an arithmetic
	circuit $\cC(\bx, \by)$ of size $s = {\rm poly}(n)$, where
	$\bx = (x_{ijk})_{i,j,k=1}^n$ is the set of variables corresponding
	to $V$ and $\by = (y_1, \ldots, y_r)$ is the set of auxiliary
	variables, with $r = {\rm poly}(n)$.
	Moreover, from the definition of the size of a succinct encoding we also have that the
	constants used in the computation of $\cC(\bx, \by)$ are rational numbers with bit
	complexity bounded by $b = {\rm poly}(n)$. In particular, $\cC(\bx, \by) \in \Q[\bx,\by]$.

	Let us consider the circuit $\cC(\bx, \by)$ as a circuit whose constants
	are in $\Q[\by]$ and whose variables are only the $\bx$ variables,
	that is, a circuit in $\Q[\by][\bx]$.
  Then, Proposition~\ref{prop:hom-components}
	tells us that there exists a homogeneous circuit $\cC_n(\bx, \by)$,
	in the $\bx$ variables, of degree $n$ and size $O(n^2 s)$ that computes the homogeneous component of $\cC(\bx,\by)$ of degree $n$ as a function of~$\bx$.
  Moreover, the constants of this circuit are a subset of the constants used in the circuit $\cC(\bx, \by)$.
  Since we consider the latter as a circuit in only the $\bx$ variables, the constants in this case are given by the elements of $\Q$ used in the computation of
  $\cC$ as well as the auxiliary variables $\by$.
  In particular, $\cC_n(\bx, \by)$ can be written in the following way:
	\begin{equation}\label{eq:deg-n-cks-1}
		\cC_n(\bx, \by) = \sum_{\bbm \in \cN_n(\bx)} f_\bbm(\by) \cdot \bbm,
	\end{equation}
	where $\cN_n(\bx)$ is the set of all monomials of degree $n$ in the
	variables $\bx$ and $f_\bbm(\by)$ are polynomials in the variables $\by$ of
	degree at most~$2^s$, as the circuit $\cC$ has size at most $s$.

	In Proposition~\ref{prop:least-degree-torus} below, we will show that the invariants
	of minimum degree of our action are in degree $n$, and these
	are spanned by the (maximum) $3$-dimensional matching monomials.
  Thus, if a monomial of degree $n$ is invariant under our action, it must be the case that this monomial corresponds to a $3$-dimensional matching.
  Moreover, the action maps any monomial (invariant or not) to a constant times itself. 
  As $\cC_n(\bx, \by)$ must only compute invariant polynomials, this
	implies that equation~\eqref{eq:deg-n-cks-1} is actually of the following form:
	\begin{equation}\label{eq:deg-n-cks-2}
		\cC_n(\bx, \by) = \sum_{\bbm \in \cM_n(\bx)}
		f_\bbm(\by) \cdot \bbm,
	\end{equation}
	where $\cM_n(\bx)$ is the set of all $3$-dimensional matching monomials
	in the variables $\bx$.
  Moverover, since $\cC(\bx,\by)$ succinctly encodes of a set of generators, the span of $\{\cC_n(\bx,\alpha)\}_{\alpha\in\C^r}$ must necessarily be the same as the span of the 3-dimensional matching polynomials.

	We will now show that the existence of the circuit $\cC_n(\bx, \by)$ implies that
	${\rm NP\subseteq P/poly}$. For that purpose, we will show that given $\cC_n(\bx, \by)$ one can
	solve the $3$-dimensional matching problem in ${\rm P/poly}$.
	Let $H$ be a tripartite 3-uniform hypergraph, whose edges are given by a subset
	$E \subseteq [n] \times [n] \times [n]$.
	We can associate to this graph the tensor $v \in V$ where
	$v_{ijk} = 1$ if hyperedge $(i,j,k) \in E$ and $v_{ijk} = 0$ otherwise.
	Note that $H$ has a $3$-dimensional matching of size $n$ if and only if at least one
	of the $3$-dimensional matching monomials \emph{does not vanish} on our tensor $v$.
	This last
	condition is equivalent to the fact that the circuit $\cC_n(v, \by)$ does not compute the
	zero polynomial (as we know that the span of the set $\{\cC_n(\bx, \alpha)\}_{\alpha \in \C^r}$
	is the same as the span of all $3$-dimensional matching monomials).
	Thus, to solve the $3$-dimensional
	matching problem in ${\rm P/poly}$ it is enough to give a randomized polynomial time algorithm
	for testing whether $\cC_n(v, \by)$ is the zero polynomial or not.\footnote{It is enough to
	give a randomized polynomial time algorithm because we know that ${\rm BPP/poly = P/poly}$.}

	Since $\cC_n(v, \by)$ is a circuit of size ${\rm poly}(n)$ with rational constants of bit
	complexity ${\rm poly}(n)$, it computes a polynomial $P(\by)$ with rational coefficients
	having bit complexity at most $\exp(\poly(n))$ and degree at most $\exp(\poly(n))$.
  This is the setting in which
	Theorem~\ref{thm:PIT-BPP} applies, giving us the desired randomized polynomial
	time algorithm. This concludes our proof modulo Proposition~\ref{prop:least-degree-torus},
	which we will now turn our attention to.
\end{proof}



%
%
%
%

%
%
%
%
%
%

%
%
%
%
%
%
%
%

\noindent
In the following proposition, we denote by $\mathcal{S}_n$ the symmetric group on $n$ letters.


\begin{proposition}\label{prop:least-degree-torus}
	The maximum $3$-dimensional matching monomials  $\prod_{i=1}^n x_{i \sigma(i) \tau(i)}$,
	where $\sigma, \tau \in \mathcal{S}_n$, span the invariants of degree
	$n$ of the natural action of $G = \ST_n(\C) \times \ST_n(\C) \times \ST_n(\C)$
	on $V = \C^n \otimes \C^n \otimes \C^n$. Moreover, there are no nonconstant invariants of degree less than $n$ for this action.
\end{proposition}

\begin{proof}
  Since the action maps any monomial to a constant times itself, it is easy to see that the invariant polynomials are generated by invariant monomials.
	To prove the proposition, it is therefore enough to show that the matching monomials are invariant, that there are no other invariant monomials of degree~$n$, and that there are no invariant monomials of smaller degree.

  We first prove that the matching monomials are invariant.
	Note that the natural action of $G$ on $V$ induces the following action on the variables $x_{ijk}$:
	$(\ba, \bb, \bc) \cdot x_{ijk} = (a_i b_j c_k)^{-1} \cdot x_{ijk}$.
  \footnote{The inverse comes from the general formula $(g\cdot p)(v) := p(g^{-1} \cdot v)$. }
  Additionally, note that
	$\prod_{\ell=1}^n a_\ell = \prod_{\ell=1}^n b_\ell = \prod_{\ell=1}^n c_\ell = 1$. Given a
	matching monomial $\prod_{i =1}^n x_{i \sigma(i) \tau(i)}$, we therefore have that
	\begin{align*}
		(\ba, \bb, \bc) \cdot \prod_{i =1}^n x_{i \sigma(i) \tau(i)}
		& = \prod_{i =1}^n \left( (a_i b_{\sigma(i)} c_{\tau(i))})^{-1} \cdot x_{i \sigma(i) \tau(i)}  \right) \\
		& = \prod_{i=1}^n (a_i b_{\sigma(i)} c_{\tau(i)})^{-1}  \cdot
		\prod_{i =1}^n x_{i \sigma(i) \tau(i)} \\
		& = \prod_{i =1}^n x_{i \sigma(i) \tau(i)}
	\end{align*}
%
	where in the last equality we note that for any permutation $\sigma \in \mathcal{S}_n$ (or $\tau$)
	we have $1= \prod_{\ell=1}^n a_\ell = \prod_{\ell=1}^n a_{\sigma(\ell)}$ (and similarly for
	$\bb$ and $\bc$). This proves that the matching monomials are invariant monomials of the natural
	$G$-action on $V$.

	Now, let us prove that no other monomial of degree $n$ is an invariant for this action.
	Let $\prod_{m=1}^n x_{i_m j_m k_m}$ be a monomial, where $(i_m, j_m, k_m) \in [n]^3$, that is not a matching monomial.
	Then there exists some coordinate, w.l.o.g.\ the first coordinate, for which the set~$\{i_m\}_{m=1}^n$ is a strict subset of~$[n]$.
  Equivalently, there is an element $\ell \in [n]$ such that $\ell \not\in \{i_m\}_{m=1}^n$.
  W.l.o.g., we can assume that
	$\ell = 1$. Thus, the action of $\ba = (\alpha^{n-1}, \alpha^{-1}, \ldots, \alpha^{-1}),
	\bb = \bc = (1, \ldots, 1)$  on our monomial $\prod_{m=1}^n x_{i_m j_m k_m}$ is as follows:
	\begin{equation*}
		(\ba, \bb, \bc) \cdot \prod_{m=1}^n x_{i_m j_m k_m}
		= \alpha^n \cdot \prod_{m=1}^n x_{i_m j_m k_m}
	\end{equation*}
	which proves that this monomial is not an invariant.
  This completes the proof that the matching
	monomials span the invariants of degree $n$.

	Now we are left with proving that there are no nonconstant monomials of degree less than~$n$ that are invariant.
  Note that if we have a monomial with degree less than $n$, we can represent it as
	$\prod_{m=1}^d x_{i_m j_m k_m}$, where $d < n$ and by the pigeonhole principle, we know that
	there exists $\ell \in [n]$ such that $\ell$ does not appear as a first coordinate entry in
	the set of tuples $\{(i_m, j_m, k_m)\}$. If $d>0$ then, analogously to the previous paragraph, we know
	that such monomials cannot be invariants of the natural action of $G$ over $V$, therefore showing
	that no nonconstant monomial of degree $< n$ can be an invariant. This completes the proof.
\end{proof}






\section{Invariant Theory for \texorpdfstring{$\SL_n(\C)$}{SL\_n(C)} and Mulmuley's conjecture} \label{Sec:Ketan}

%

In this section, we will give another example of a group action on tensors for which any set of generating invariants is hard to compute, i.e., we will prove Theorem~\ref{MC-cex}. Even though the previous section already gives a counterexample, this example illustrates something more. The feature of this group action is that invariants of minimial degree span a $1$-dimensional space. In other words, up to scaling, we have a unique invariant of minimal degree. This unique invariant in the minimal degree is called the \emph{hyperpfaffian polynomial} (introduced by Barvinok in 1995 as a natural generalization of the well-known Pfaffian polynomial to higher order tensors). We then study the hyperpfaffian's computational complexity and prove that it is VNP-complete. The importance of this example is that such a unique invariant in the minimal degree is {\em essential} in any generating set. So, it is not even possible to give a generating set consisting of invariant polynomials that are easy to compute, even if we remove all restrictions on the size of the generating set.
Moreover, the group action is by $\SL_n(\C)$ rather than a torus.
Therefore our counterexample disproves Mulmuley's original conjecture in a strong sense.
%

%
%
%


%
\subsection{Invariant Rings and Symmetric Tensors}\label{sec:invariants}
%
The special linear group $\SL_n(\C)$ consists of complex $n\times n$-matrices with unit determinant and acts canonically on $\C^n$ by matrix-vector multiplication.
This action extends to any $m$-th tensor power $\tensor^m \C^n$ by
\begin{equation}\label{eq:SL on tensor power}
g \cdot (v_1 \otimes \cdots \otimes v_m) := (g \cdot v_1) \otimes \cdots \otimes (g \cdot v_m)
\end{equation}
and linear continuation.
We will always use the standard bilinear form on $\tensor^m \C^n$ that satisfies $\langle g^T \cdot v, w\rangle = \langle v, g \cdot w \rangle$ for all $v,w \in \tensor^m\C^n$, $g \in \SL_n(\C)$.
%

Let~$V$ be an arbitrary finite-dimensional $\SL_n(\C)$-representation (such as $V = \tensor^m\C^n$).
Then $\SL_n(\C)$ also acts on~$\C[V]_d$, the vector space of degree-$d$ homogeneous polynomials on~$V$, via the formula
\[
(g \cdot p)(v) := p(g^T \cdot v),
\]
where $p \in \C[V]_d, g \in G$ and $v \in V$.
The formula above is the dual representation of the action on the ring of all polynomial functions~$\C[V] = \bigoplus_{d=0}^\infty \C[V]_d$ that we explained in in Section 1.2 of the original publication.
Using the dual here is only for presentation purposes, as it gives a clearer connection to multilinear algebra as follows.
%
Note that a polynomial $p \in \C[V]$ is invariant if and only if $\forall g \in \SL_n(\C)$ we have $g \cdot p=p$.
%
%

It is convenient to identify polynomial functions with symmetric tensors.
Note that $\SL_n(\C)$ acts canonically on any $d$-th tensor power $\tensor^d V$ of~$V$.
This action restricts to the $d$-th symmetric tensor power~$\Sym^d V$, i.e., the $\mathcal{S}_d$-invariant subspace of $\tensor^d V$. Recall that $\mathcal{S}_d$ is the symmetric group on $d$ letters; it acts on $V^{\otimes d}$ by permuting tensor factors.
For any~$t \in \Sym^d V$, we can define a homogeneous degree-$d$ polynomial~$p\in\C[V]_d$ by~$p(v) := \langle t, v^{\otimes d}\rangle$.
Here we use the quadratic form on $\Sym^d V$ induced by a non-degenerate bilinear form on $V$ that satisfies $\langle g^t\cdot v, w\rangle = \langle v,g w\rangle$ for all $v,w \in V$, $g \in \SL_n(\C)$ as above.
Then, $p$ is invariant if and only if the symmetric tensor~$t$ is invariant, i.e., if $\forall g \in \SL_n(\C)$ we have~$g \cdot t=t$.
We will tacitly go back and forth between symmetric tensors in $\Sym^d \tensor^m \C^n$ and homogeneous polynomials in~$\C[\tensor^m \C^n]_d$.

Now, we turn to studying hyperpfaffians.





%
%
%
%

%
%
%
%
%
%
%

%
\subsection{Hyperpfaffians}
%
The \emph{Pfaffian} is the unique (up to scale) homogeneous $\SL_{2n}(\C)$-invariant of degree $n$ on~$\C^{2n}\otimes\C^{2n}$.
There are no $\SL_{2n}(\C)$-invariants in lower degrees.
%
If we identify $\C^{2n} \otimes \C^{2n}$ with the space of complex $2n\times 2n$ matrices~$A$, then the Pfaffian is invariant under the action of $\SL_{2n}(\C)$ given by $g \cdot A := gAg^T$.
The defining property of the Pfaffian generalizes to tensors of even order as follows (the classical Pfaffian is the special case of~$k=1$):

\begin{proposition}\label{pro:nolowerdegreeinv}
For any $k$ and $n$, there is a unique (up to scale) homogeneous $\SL_{2kn}(\C)$-in\-va\-ri\-ant polynomial~$\Pf_{k,n}$ of degree~$n$ on $\tensor^{2k} \C^{2kn}$.
$\Pf_{k,n}$ identifies with the symmetric tensor $e_1 \wedge \cdots \wedge e_{2kn} \in \Sym^n \tensor^{2k}\C^{2kn}$.
There are no nonconstant $\SL_{2kn}(\C)$-invariants of lower degree.
\end{proposition}

\noindent
Before proving Proposition~\ref{pro:nolowerdegreeinv} we recall some representation theory.
The material is well-known, and we refer to standard texts (e.g., \cite{FH}) for details.
A \emph{partition} $\la$ is a nonincreasing sequence of natural numbers with finite support.
We write $\la \vdash_n m$ to say that~$|\la| := \sum_i \la_i = m$ and $\la_{n+1}=0$.
If $\la_{n+1}=0$, then we say that $\la$ is an $n$-partition.
The irreducible polynomial $\GL_n(\C)$-representations are indexed by $n$-partitions.
For a partition~$\la$, let~$\{\la\}$ denote the irreducible $\GL_n(\C)$-representation corresponding to~$\la$.
Restricted to $\SL_n(\C)$, the representation $\{\la\}$ is trivial if and only if $\la_1=\ldots=\la_n$; note that this implies that $n \mid m$.
The irreducible representations of $\mathcal{S}_m$ are indexed by partitions~$\la$ with~$|\la|=m$.
Let $[\la]$ denote the irreducible $\mathcal{S}_m$-representation corresponding to~$\la$.
%

Consider $\tensor^m \C^n$. This space has an action of $\SL_n(\C)$ by~\eqref{eq:SL on tensor power}, but also an action of $\mathcal{S}_m$ that permutes the tensor factors.
Both actions commute, so we have an action of the product group $\SL_n(\C)\times\mathcal S_m$. The following well-known result will be crucial for our purposes.

\begin{theorem} [Schur--Weyl duality]
As an $\SL_n(\C) \times \mathcal{S}_m$-representation, we have the decomposition:
\[ 
\tensor^m \C^n = \bigoplus_{\la \vdash_n m} \{\la\}\otimes[\la].
\] 
\end{theorem}

\noindent
Using Schur--Weyl duality, one sees immediately that $\tensor^m \C^n$ contains nonzero $\SL_n(\C)$-invariant vectors if and only if $n \mid m$.
This is because a vector is invariant if and only if it spans a trivial irreducible representation -- but $\{\la\}$ is trivial if and only if $\la_1=\ldots=\la_n$, as mentioned above.
For $m=n$, there is a unique (up to scale) $\SL_n(\C)$-invariant vector.
This is because the invariants in $\tensor^n \C^n$ correspond to the component $\{1^n\} \otimes [1^n]$, where we write~$1^n$ for the partition~$\la_1=\ldots=\la_n=1$.
Here, $\{1^n\}$ is the trivial representation of $\SL_n(\C)$ and $[1^n]$ is one-dimensional, as it is the sign representation of~$\mathcal{S}_n$.
Thus the space of invariants is one-dimensional. This unique vector (up to scale) is given by the wedge product $e_1 \wedge e_2 \wedge \cdots \wedge e_n$, where $a \wedge b := \frac 1 2 (a \otimes b - b \otimes a)$, and higher order wedge products are defined analogously.

\begin{proof}[Proof of Proposition~\ref{pro:nolowerdegreeinv}]
It suffices to show that $\Sym^d \tensor^{2k} \C^{2kn}$ contains no $\SL_{2kn}(\C)$-invariant vector if $0<d<n$ and that it contains a unique such vector if $d=n$.
Note that $\Sym^d \tensor^{2k} \C^{2kn}$ is a subspace of $\tensor^d \tensor^{2k} \C^{2kn} \simeq \tensor^{2kd} \C^{2kn}$.
Thus the first claim holds since $\tensor^{2kd} \C^{2kn}$ contains $\SL_{2kn}(\C)$-invariant vectors only if $2kn \mid 2kd$. Thus if $0 < d < n$, there are no invariants.
For~$d=n$, $\tensor^d \tensor^{2k} \C^{2kn} \simeq \tensor^{2kn} \C^{2kn}$ contains the unique $\SL_{2kn}(\C)$-invariant vector $v = (e_1 \wedge \cdots \wedge e_{2k}) \wedge \cdots \wedge (e_{2k(n-1)+1} \wedge \cdots \wedge e_{2kn})$.
It remains to show that $v$ is symmetric, i.e., an element of~$\Sym^d \tensor^{2k} \C^{2kn}$, which is a subspace of $\tensor^d \tensor^{2k} \C^{2kn}$.
But this is easy to see since each of the~$d$ blocks has even size~$2k$ and the wedge product is skew-commutative.
This proves the second claim.
\end{proof}

%
%

%
%

%
%

%
%
%
%
%
%
%
%
%
%
%
%
%

%
%
%
%
%
%

\noindent
The polynomial~$\Pf_{k,n}$ was introduced in~\cite[Def.~3.4]{Bar:95} in its monomial presentation, where it is called the \emph{hyperpfaffian}.
Note that, for fixed $k$, $\Pf_k := (\Pf_{k,1},\Pf_{k,2},\dots)$ is a p-family (i.e., both the degree and the number of variables are polynomially bounded), since~$\Pf_{k,n}$ has degree~$n$ and $(2kn)^{2k}$ variables.
The monomial presentation in~\cite{Bar:95} immediately yields that $\Pf_k \in \VNP$.
\begin{theorem}
For even $k$, $\Pf_k$ is \VNP-complete.
\end{theorem}
\begin{proof}
We present a projection of $\Pf_{k,d}$ to the $d\times d$ permanent.
The same projection yields the determinant if $k$ is odd, which explains why the proof does not work for the classical Pfaffian ($k=1$).
The case~$k=2$ is enough to disprove Mulmuley's conjecture.

By Proposition~\ref{pro:nolowerdegreeinv}, the Pfaffian $\Pf_{k,d}$ identifies with the symmetric tensor
\begin{align*}
  v := e_1 \wedge \cdots \wedge e_{2kd} \in \Sym^d \tensor^{2k}\C^{2kd}.
\end{align*}
%
Thus, the evaluation~$\Pf_{k,d}(p)$ at a tensor $p \in \tensor^{2k} \C^{2kd}$ is given by $\langle v, p^{\otimes d}\rangle$ (cf.~\cite[Sec.~4.2(A)]{Ike:12}).
We choose
\[
\begin{aligned}
p ={}& \sum_{i,j=0}^{d-1} x_{i+1,j+1} (e_{1+2ki} \otimes e_{2+2ki} \otimes \cdots \otimes e_{k+2ki}\\
&{}\otimes e_{k+1+2kj} \otimes e_{k+2+2kj} \otimes \cdots \otimes e_{2k+2kj}),
\end{aligned}
\]
where the $x_{i,j}$ ($1 \leq i,j \leq d$) are formal variables.

The point $p$ is parametrized linearly by the $x_{i,j}$, so the evaluation of $\Pf_{k,d}$ at $p$ is a projection of $\Pf_{k,d}$.
We verify that the evaluation of $\Pf_{k,n}$ at $p$ gives the $d \times d$ permanent (up to a constant nonzero scalar) as follows.

{
\fontsize{10}{12}
\selectfont
\begin{align*}
\quad p^{\otimes d} &=\hspace{-.5cm}\sum_{i_1,j_1,\ldots,i_d,j_d=0}^{d-1}\hspace{-.5cm} x_{i_1+1,j_1+1} \cdots x_{i_d+1,j_d+1}
(e_{1+2ki_1} \otimes e_{2+2ki_1} \otimes \cdots \otimes e_{k+2ki_1}\\ &\quad{}\otimes e_{k+1+2kj_1} \otimes e_{k+2+2kj_1} \otimes \cdots \otimes e_{2k+2kj_1})\\ &\quad\otimes \cdots \otimes
(e_{1+2ki_d} \otimes e_{2+2ki_d} \otimes \cdots \otimes e_{k+2ki_d}\\ &\quad{}\otimes e_{k+1+2kj_d} \otimes e_{k+2+2kj_d} \otimes \cdots \otimes e_{2k+2kj_d})
\end{align*}
}
and by linearity
{
\fontsize{10}{12}
\selectfont
\begin{align*}
\quad \langle v,p^{\otimes d}\rangle &= \hspace{-.5cm}\sum_{i_1,j_1,\ldots,i_d,j_d=0}^{d-1}\hspace{-.5cm} x_{i_1+1,j_1+1} \cdots x_{i_d+1,j_d+1}
\langle v,(e_{1+2ki_1} \otimes e_{2+2ki_1} \otimes \cdots \otimes e_{k+2ki_1}\\ &\quad{}\otimes e_{k+1+2kj_1} \otimes e_{k+2+2kj_1} \otimes \cdots \otimes e_{2k+2kj_1})\\ &\quad\otimes \cdots \otimes
(e_{1+2ki_d} \otimes e_{2+2ki_d} \otimes \cdots \otimes e_{k+2ki_d}\\ &\quad{}\otimes e_{k+1+2kj_d} \otimes e_{k+2+2kj_d} \otimes \cdots \otimes e_{2k+2kj_d})\rangle
\end{align*}
}
A crucial property of $v$ is that $\langle v,e_{\pi(1)} \otimes e_{\pi(2)} \otimes \cdots \otimes e_{\pi(n)}\rangle \neq 0$ iff $\pi$ is a permutation, in which case it is equal to the sign of the permutation.
It follows that the nonzero summands in $\langle v,p^{\otimes d}\rangle$ are precisely those for which~$i=(i_1,\ldots,i_d)$ and $j=(j_1,\ldots,j_d)$ are permutations of $\{0,\ldots,d-1\}$.
For a single summand with $i$ and $j$ permutations we see:
\begin{align*}
&\quad x_{i_1+1,j_1+1} \cdots x_{i_d+1,j_d+1}
\langle v,(e_{1+2ki_1} \otimes e_{2+2ki_1} \otimes \cdots \otimes e_{k+2ki_1} \otimes e_{k+1+2kj_1} \otimes e_{k+2+2kj_1} \otimes \cdots \otimes e_{2k+2kj_1})\\ &\quad\quad\quad\quad\quad\quad\quad\quad\otimes \cdots \otimes
(e_{1+2ki_d} \otimes e_{2+2ki_d} \otimes \cdots \otimes e_{k+2ki_d} \otimes e_{k+1+2kj_d} \otimes e_{k+2+2kj_d} \otimes \cdots \otimes e_{2k+2kj_d})\rangle
\\&= \sgn(i)^k \sgn(j)^k x_{i_1+1,j_1+1} \cdots x_{i_d+1,j_d+1}.
\end{align*}
Hence, for even $k$ we obtain $\langle v,p^{\otimes d}\rangle = d! \per_d$.
\end{proof}

\noindent
Finally, we put together the preceding results to prove Theorem~\ref{MC-cex}.

\begin{theorem}[Theorem~\ref{MC-cex}, restated]
Let $k \geq 2$ be even.
Consider the action of $G = \SL_{2kn}(\C)$ on $V = \tensor^{2k} \C^{2kn}$. Then any set of generators for the invariant ring cannot have a polynomial sized (in $n$) succinct encoding, unless ${\rm VP = VNP}$.
\end{theorem}
\begin{proof}

We summarize the results so far. Let $k \geq 2$. Consider the action of $G = \SL_{2kn}(\C)$ on $V = \tensor^{2k} \C^{2kn}$. Then:

\begin{enumerate}
\item There are no homogeneous invariant polynomials of degree $< n$.

\item The space of homogeneous invariant polynomials of degree $n$ is $1$-dimensional, and spanned by the hyperpfaffian polynomial $\Pf_{k,n}$.

\item The hyperpfaffian polynomial $\Pf_{k,n}$ is VNP-complete.
\end{enumerate}
The rest of the proof proceeds along the same lines as in the proof of Theorem~\ref{thm:torus-action} in Section~\ref{sec:torus}. If we had a poly-sized succinct encoding for the generators of this invariant ring, then one would be able to extract the lowest degree part, which would yield a poly-sized circuit computing $\Pf_{k,n}$.
This is not possible unless ${\rm VP = VNP}$, since $\Pf_{k,n}$ is VNP-complete.
\end{proof}




\begin{thebibliography}{99}
\bibitem{AB03}
Manindra Agrawal and Somenath Biswas.
\newblock Primality and identity testing via chinese remaindering.
\newblock {\em Journal of the ACM (JACM)}, 50(4):429--443, 2003.

\bibitem{Bar:95}
Alexander~I. Barvinok.
\newblock New algorithms for linear $k$-matroid intersection and matroid
  $k$-parity problems.
\newblock {\em Mathematical Programming}, 69(1):449--470, 1995.

\bibitem{DK}
Harm Derksen and Gregor Kemper.
\newblock {\em Computational invariant theory}, volume 130 of {\em
  Encyclopaedia of Mathematical Sciences}.
\newblock Springer, Heidelberg, enlarged edition, 2015.
\newblock With two appendices by Vladimir L. Popov, and an addendum by Norbert
  A'Campo and Popov, Invariant Theory and Algebraic Transformation Groups,
  VIII.

\bibitem{DM-exp}
Harm {Derksen} and Visu {Makam}.
\newblock {An exponential lower bound for the degrees of invariants of cubic
  forms and tensor actions}.
\newblock {\em arXiv e-prints}, page arXiv:1902.10773, Feb 2019.

\bibitem{FH}
William Fulton and Joe Harris.
\newblock {\em Representation theory}, volume 129 of {\em Graduate Texts in
  Mathematics}.
\newblock Springer-Verlag, New York, 1991.
\newblock A first course, Readings in Mathematics.

\bibitem{Ike:12}
Christian Ikenmeyer.
\newblock {\em Geometric Complexity Theory, Tensor Rank, and
  {L}ittlewood-{R}ichardson Coefficients}.
\newblock PhD thesis, Institute of Mathematics, University of Paderborn, 2012.
\newblock Online available at
  \url{http://digital.ub.uni-paderborn.de/ubpb/urn/urn:nbn:de:hbz:466:2-10472}.

\bibitem{K72}
Richard~M Karp.
\newblock Reducibility among combinatorial problems.
\newblock In {\em Complexity of computer computations}, pages 85--103.
  Springer, 1972.

\bibitem{S79}
Jacob~T Schwartz.
\newblock Probabilistic algorithms for verification of polynomial identities.
\newblock In {\em International Symposium on Symbolic and Algebraic
  Manipulation}, pages 200--215. Springer, 1979.

\bibitem{S73}
Volker Strassen.
\newblock Vermeidung von divisionen.
\newblock {\em Journal f{\"u}r die reine und angewandte Mathematik},
  264:184--202, 1973.

\bibitem{SY10}
Amir Shpilka, Amir Yehudayoff, et~al.
\newblock Arithmetic circuits: A survey of recent results and open questions.
\newblock {\em Foundations and Trends{\textregistered} in Theoretical Computer
  Science}, 5(3--4):207--388, 2010.

\bibitem{Wehlau}
David~L. Wehlau.
\newblock Constructive invariant theory for tori.
\newblock {\em Ann. Inst. Fourier (Grenoble)}, 43(4):1055--1066, 1993.

\end{thebibliography}
\end{document}
